\documentclass{article}

% Math
\usepackage{amsmath, amssymb, amsthm}

% Tables
\usepackage{booktabs}
\usepackage{tabularx}

% Figures
\usepackage{graphicx}
\usepackage{subcaption}

% Commands
\newcommand{\algorithmicbreak}{\textbf{break}}
\newcommand{\BREAK}{\STATE \algorithmicbreak}

% Misc
\usepackage{enumerate}
\usepackage{hyperref}
\usepackage{float}
\usepackage{algorithm}
\usepackage{algpseudocode}
\usepackage{biblatex}

\title{Econ 527 HW3}
\author{Camille Bergeron}
\date{May 2026}

\begin{document}
\maketitle

% ============================================================
\section{The Hugget Model}
% ============================================================
From calculation, the average income for the borrowing limit i.e. the minimum value on the welath grid is around 1.37. Thus, the maximum is 65.2287.

The equilibrium interest rate $r^*$ is detemrined when 
\begin{equation*}
    A(r^*)=0
\end{equation*}

Thus, we will use the bisection algorithm find $r^*$, setting the lower bound $r_l=0$ amnd  

\end{document}